\subsection{TYPICAL QUANTIFIERS}

Let \(\pmb{A}\) and \(\pmb{R}\) be assemblies and let \(\pmb{x}\) be a letter. We denote the assembly \((\exists x)(\pmb{A}\) and \(\pmb{R})\) by \((\exists_{\pmb{A}}x)\pmb{R}\) , and the assembly

\[
\text {   ``not   } (\exists_ {A} x) \text {   (not   } R)
\]

by \((\forall_{A}x)R\) . The abbreviating symbols \(\exists_{A}\) and \(\forall_{A}\) are called typical quantifiers. Observe that the letter x does not appear in the assemblies denoted by \((\exists_{A}x)R\) , \((\forall_{A}x)R\) .

CS10. Let A and R be assemblies and let x and \(x'\) be letters. If \(x'\) appears neither in R nor in A, then \((\exists_{A}x)R\) and \((\forall_{A}x)R\) are respectively identical with \((\exists_{A'}x')R'\) and \((\forall_{A'}x')R'\) , where \(R'\) is \((x'|x)R\) and \(A'\) is \((x'|x)A\) .

CS11. Let \(A, R, U\) be assemblies, and let \(x, y\) be distinct letters. If \(x\) does not appear in \(U\) , the assemblies \((U|y)(\exists_A x)R\) and \((U|y)(\forall_A x)R\) are respectively identical with \((\exists_{A'}x)R'\) and \((\forall_{A'}x)R'\) , where \(R'\) is \((U|y)R\) and \(A'\) is \((U|y)A\) .

These rules are immediate consequences of criteria CS8, CS9 (no. 1), CS5 (§1, no. 2), and CS6 (§3, no. 4).

CF12. Let A and R be relations in G, and let x be a letter. Then

\[
(\exists_ {A} x) R \quad \text{and} \quad (\forall_ {A} x) R
\]

are relations in \(\mathfrak{C}\) .

This follows directly from CF11 (no. 1), CF9 (§3, no. 4), and CF2 (§1, no. 4).

Intuitively, consider A and R as expressing properties of x. It may happen that in a series of proofs, we are concerned only with objects satisfying A. To say that there exists an object satisfying A such that R means that there exists an object such that ``A and R''; whence the definition of \(\exists_{A}\) . To say that all objects which satisfy \(A\) have the property \(R\) means that there is no object satisfying \(A\) such that ``not \(R\)''; whence the definition of \(\forall_{A}\) . In practice, these signs are replaced by various phrases, depending on the nature of the relation \(A\) . * For example: ``for all integers \(x\) , \(R\)''; ``there exists an element \(x\) of the set E such that \(R\)''; and so on. *

C35. Let A and R be relations in G, and let x be a letter. Then the relations \((\forall_{A}x)R\) and \((\forall x)(A\Longrightarrow R)\) are equivalent in G.

For the relation \((\forall_{A}x)R\) is identical with

\[
\text{``not } (\exists x) (A \text{ and } (\text{not } R))\text{''}.
\]

Now, ``A and (not R)'' is equivalent in \(C_{0}\) to ``not \((A \Longrightarrow R)\) ''; therefore ``not \((\exists x)(A \text{ and (not } R))\)'' is equivalent in \(C_{0}\) to

\[
\text{``not } (\exists x) (\text{not } (A \Rightarrow R))\text{''},
\]

by C31 (no. 3), and the latter relation is identical with \((\forall x)(A \Rightarrow R)\) . The criterion is therefore established in \(\mathfrak{G}_0\) , and consequently in \(\mathfrak{G}\) .

We shall often have to prove relations of the form \((\forall_{A}x)R\) ; generally we shall use one of the following two criteria:

C36. Let A and R be relations in C, and let x be a letter. Let \(C'\) be the theory obtained by adjoining A to the axioms of C. If x is not a constant of C, and if R is a theorem in \(C'\) , then \((\forall_{A}x)R\) is a theorem in C.

For \(A \Longrightarrow R\) is a theorem in \(\mathfrak{C}\) , by the criterion of deduction, therefore \((\forall_A x) R\) is a theorem in \(\mathfrak{C}\) by C27 (no. 1) and C35.

In practice, we indicate that we are going to use this rule by a phrase such as ``Let \(\pmb{x}\) be any element such that \(A\)''. In the theory \(\mathfrak{C}'\) so defined, we seek to prove \(R\) . Of course, we cannot assert that \(R\) itself is a theorem in \(\mathfrak{C}\) .

C37. Let A and R be relations in C and let x be a letter. Let \(C'\) be the theory obtained by adjoining the relations A and ``not R'' to the axioms of C. If x is not a constant of C, and if \(C'\) is contradictory, then \((\forall_{A}x)R\) is a theorem in C.

For the theory \(\mathfrak{C}'\) is equivalent to the theory obtained by adjoining ``not \((A \Longrightarrow R)\)'' to the axioms of \(\mathfrak{C}\) . By the method of reductio ad absurdum, \(A \Longrightarrow R\) is a theorem in \(\mathfrak{C}\) , and therefore so is \((\forall_A x)R\) by C27 (no. 1) and C35.

In practice we say: ``Suppose that there exists an object \(x\) satisfying \(A\) for which \(R\) is false'', and seek to establish a contradiction.

The properties of typical quantifiers are analogous to those of quantifiers :

C38. Let A and R be relations in G, and let x be a letter. Then the relations

\[
\begin{array}{l} \text { not } (\forall_ {A} x) R \Longleftrightarrow (\exists_ {A} x) (\text { not } R), \\ \text { not } (\exists_ {A} x) R \Longleftrightarrow (\forall_ {A} x) (\text { not } R) \end{array}
\]

are theorems in \(\mathfrak{C}\) .

C39. Let A, R, and S be relations in C, and let x be a letter which is not a constant of C. If the relation \(A \Longrightarrow (R \Longrightarrow S)\) {[}resp. \(A \Longrightarrow (R \Longleftrightarrow S)\) {]} is a theorem in C, then the relations

\[
\begin{array}{c c} (\exists_ {A} x) R \Longrightarrow (\exists_ {A} x) S, & (\forall_ {A} x) R \Longrightarrow (\forall_ {A} x) S \\ \text{[resp. } (\exists_ {A} x) R \Longleftrightarrow (\exists_ {A} x) S, & (\forall_ {A} x) R \Longleftrightarrow (\forall_ {A} x) S \text{]} \end{array}
\]

are theorems in \(\mathfrak{C}\) .

C40. Let A, R, and S be relations in C and let x be a letter. Then the relations

\[
\begin{array}{c} (\forall_ {A} x) (R \text {   and   } S) \iff ((\forall_ {A} x) R \text {   and   } (\forall_ {A} x) S), \\ (\exists_ {A} x) (R \text {   or   } S) \iff ((\exists_ {A} x) R \text {   or   } (\exists_ {A} x) S) \end{array}
\]

are theorems in \(\mathfrak{C}\) .

C41. Let A, R, and S be relations in C, and let x be a letter which does not appear in R. Then the relations

\[
\begin{array}{l} (\forall_ {A} x) (R \text {   or   } S) \iff (R \text {   or   } (\forall_ {A} x) S), \\ (\exists_ {A} x) (R \text {   and   } S) \iff (R \text {   and   } (\exists_ {A} x) S) \end{array}
\]

are theorems in \(\mathfrak{C}\) .

C42. Let A, B, R be relations in C and let x and y be letters. If x does not appear in B, and if y does not appear in A, then the relations

\[
\begin{array}{l} (\forall_ {A} x) (\forall_ {B} y) R \iff (\forall_ {B} y) (\forall_ {A} x) R, \\ (\exists_ {A} x) (\exists_ {B} y) R \iff (\exists_ {B} y) (\exists_ {A} x) R, \\ (\exists_ {A} x) (\forall_ {B} y) R \implies (\forall_ {B} y) (\exists_ {A} x) R \end{array}
\]

are theorems in \(\mathfrak{C}\) .

By way of example, let us prove part of C42. The relation

\[
(\exists_ {A} x) (\exists_ {B} y) R
\]

is identical with \((\exists x)(A\) and \((\exists y)(B\) and \(R)\) , and therefore (because \(y\) does not appear in \(A\) ) is equivalent in \(\mathcal{T}_0\) to

\[
(\exists x) (\exists y) (A \text {   and   } (B \text {   and   } R))
\]

by C33 and C31. Likewise, \((\exists_{B}x)(\exists_{A}y)R\) is equivalent to

\((\exists y)(\exists x)(B\) and \((A\) and \(R)\) .

Now apply C31 and C34 (no. 3).

* As an example of the application of these criteria, consider the following relation: ``the sequence of real-valued functions \((f_{n})\) converges uniformly to 0 in {[}0, 1{]}''. This means ``for each \(\varepsilon > 0\) there exists an integer n such that for each \(x \in [0, 1]\) and each integer \(m \geqslant n\) we have \(|f_{m}(x)| \leqslant \varepsilon\) ''. Suppose we wish to take the negation of this relation (for example, to obtain a proof by contradiction); the criterion C38 shows that this negation is equivalent to the following relation: ``there exists an \(\varepsilon > 0\) such that for each integer n there exists an \(x \in [0, 1]\) and an \(m \geqslant n\) for which \(|f_{m}(x)| > \varepsilon\) ''.

